\section*{Capítulo \ref{ch:g9:acids-bases-ph} --- Ácidos, bases e pH}

\begin{solution}{exo:g9:acids-bases-ph:1}
\omterm{def:g9:acids-bases-ph:ph}{pH} 3: ácida. \omterm{def:g9:acids-bases-ph:ph}{pH} 7: neutra. \omterm{def:g9:acids-bases-ph:ph}{pH} 11: básica.
\end{solution}

\begin{solution}{exo:g9:acids-bases-ph:2}
Ácida: \omterm{def:g9:ions:ion}{íons} \ce{H+}. Básica: \omterm{def:g9:ions:ion}{íons} \ce{OH-}.
\end{solution}

\begin{solution}{exo:g9:acids-bases-ph:3}
Suco de limão (2), café (5), água pura (7), água do mar (8.1), amoníaco
doméstico (12).
\end{solution}

\begin{solution}{exo:g9:acids-bases-ph:4}
Uma tira num pratinho limpo e seco é tocada com uma gota da solução
levada por um bastão de vidro limpo; a cor é comparada imediatamente com
a tabela.
\end{solution}

\begin{solution}{exo:g9:acids-bases-ph:5}
Cerca de 5 depois de uma diluição de dez vezes; cerca de 6 depois de uma de cem vezes.
\end{solution}

\begin{solution}{exo:g9:acids-bases-ph:6}
Não. Diluir aproxima o \omterm{def:g9:acids-bases-ph:ph}{pH} de 7, mas nunca o faz passar de 7: a própria
água é neutra, e acrescentá-la não pode fazer os \omterm{def:g9:ions:ion}{íons} \ce{OH-} ficarem
mais numerosos que os \omterm{def:g9:ions:ion}{íons} \ce{H+}.
\end{solution}

\begin{solution}{exo:g9:acids-bases-ph:7}
A \omterm{def:g6:pure-substances-and-mixtures:mixture}{mistura} libera calor. A água despejada sobre um ácido concentrado
poderia ferver na hora e espirrar ácido; o ácido despejado devagar na
água se espalha imediatamente num volume grande.
\end{solution}

\begin{solution}{exo:g9:acids-bases-ph:8}
Corrosivo: queima a pele e os olhos, ataca os metais. Usar óculos de
proteção e luvas; enxaguar qualquer respingo com muita água.
\end{solution}

\begin{solution}{exo:g9:acids-bases-ph:9}
Três diluições (de 1 para 2, de 2 para 3, de 3 para 4): $10 \times 10 \times 10 =
1000$ vezes.
\end{solution}

\begin{solution}{exo:g9:acids-bases-ph:10}
O leite de magnésia é básico (\omterm{def:g9:acids-bases-ph:ph}{pH} cerca de 10.5): ele reage com uma parte
do excesso de ácido do estômago.
\end{solution}

\begin{solution}{exo:g9:acids-bases-ph:11}
Não: com um \omterm{def:g9:acids-bases-ph:ph}{pH} de 8.1, ele continua levemente básico. Ele está indo em
direção a um \omterm{def:g9:acids-bases-ph:ph}{pH} menor, menos básico, porque o dióxido de carbono do \omterm{def:g4:air-a-mixture-of-gases:air}{ar}
\omterm{def:g2:mixing-and-dissolving:dissolve}{se dissolve} nele.
\end{solution}

\begin{solution}{exo:g9:acids-bases-ph:12}
O papel indicador é lido a olho, comparando com uma tabela de cores, com
precisão de cerca de uma unidade; o \omterm{def:g9:acids-bases-ph:ph-paper}{pHmetro} dá um número com uma ou duas
casas decimais. O \omterm{def:g9:acids-bases-ph:ph-paper}{pHmetro} é mais preciso; os dois concordam dentro da
precisão do papel.
\end{solution}

\begin{solution}{pb:g9:acids-bases-ph:1}
\textbf{1.} Ácida: ela contém principalmente \omterm{def:g9:ions:ion}{íons} \ce{H+(aq)} e
\ce{Cl-(aq)}.

\textbf{2.} O GHS05, corrosivo: um ácido ataca a pele, os olhos e os
metais.

\textbf{3.} Ele dá um número preciso, em vez de uma cor a ser comparada a
olho.

\textbf{4.} Dez vezes mais: uma unidade de \omterm{def:g9:acids-bases-ph:ph}{pH} corresponde a uma diluição
de dez vezes.

\textbf{5.} \qty{10}{L}, de \omterm{def:g9:acids-bases-ph:ph}{pH} 3.

\textbf{6.} Três: de \omterm{def:g9:acids-bases-ph:ph}{pH} 2 para 3, de 3 para 4, de 4 para 5.

\textbf{7.} $1 \times 10 \times 10 \times 10 = \qty{1000}{L}$.

\textbf{8.} Não: diluir só aproxima o \omterm{def:g9:acids-bases-ph:ph}{pH} de 7.

\textbf{9.} \ce{H+ + OH- -> H2O}.

\textbf{10.} O ácido é destruído, transformado em água, em vez de ser
espalhado num volume enorme; só é preciso um pouco de base, e não mil
litros de água.

\textbf{11.} O bicarbonato de sódio dissolvido em água (ou água com
sabão).

\textbf{12.} $1000 - 1 = \qty{999}{L}$ de água.
\end{solution}
